\clearpage
\section{Discussion}\label{section::Discussion}


This chapter provides comprehensive analysis and interpretation of the experimental results presented in Chapter 5. We investigate the fundamental reasons behind performance differences among the three Kalman state-space configurations, elucidate why the proposed physics-constrained model (Phys-5D) achieves superior performance, and analyze the synergistic interactions between system components. We further discuss theoretical implications, practical value, and limitations, providing insights for future research directions.

\subsection{Comprehensive Interpretation of Results and Core Findings}

The experimental results from Chapter 5 demonstrate the effectiveness of our detect–track–count framework. Our final system achieves state-of-the-art counting accuracy with MAPE=2.24\%, representing a significant improvement over baseline approaches. The key findings from our comprehensive evaluation are:

\textbf{Performance Summary:} The three Kalman state-space configurations show distinct performance characteristics:
\begin{itemize}
    \item CV-8D (baseline): MAPE=14.59\%, MAE=3.4, ID switches=24
    \item CA-12D (acceleration): MAPE=6.99\%, MAE=2.4, ID switches=39.15
    \item Phys-5D (physics-constrained): MAPE=2.24\%, MAE=0.75, ID switches=24.8
\end{itemize}

\textbf{Core Finding:} The results demonstrate that incorporating domain knowledge (camera geometry and physical constraints) into the motion model is far more effective than merely increasing kinematic model complexity. Phys-5D achieves the best balance between tracking accuracy and identity consistency, translating to superior counting performance.

An important observation is that application performance (counting accuracy) is not linearly correlated with traditional MOT metrics (e.g., MOTA). CA-12D and Phys-5D obtain nearly identical MOTA (67.54\% vs 67.21\%), suggesting similar per-frame detection–association accuracy. However, this obscures a critical difference: the number of IDSW is much higher for CA-12D (avg 39.15) than for Phys-5D (avg 24.80). Identity switches measure how well a tracker preserves identity consistency. For cumulative tasks like counting, which require long-term, stable, and unique identities, trajectory stability and persistent identity are more important than single-frame localization accuracy. Each identity switch can cause the same person to be mistaken for two different identities, inducing overcounts or undercounts. Hence, although CA-12D can fit instantaneous motion better, its weaker identity consistency directly leads to much higher counting errors than Phys-5D. This highlights that application-oriented evaluations must emphasize identity-consistency metrics such as IDF1 and IDSW. Phys-5D’s true advantage lies not only in accurate tracking but also in sustaining consistent identities.

\subsection{Motion Models: From Pure Kinematics to Physics-Constrained Modeling and the Trade-offs}

To uncover the origins of performance differences among motion models, we now compare the three variants evaluated—CV-8D, CA-12D, and Phys-5D—and argue why imposing physics constraints is key to success in our task.

\subsubsection{Limitations of the Baseline (CV-8D)}

As the standard baseline in MOT, the CV-8D model is our starting point. It assumes linear, constant-speed motion in the image plane. In our scenario—platform pedestrians observed from a moving train—this assumption is fundamentally violated. As the train decelerates, perspective projection causes apparent nonlinear acceleration of pedestrians in the image plane and a nonlinear enlargement of their bounding boxes.
The experimental evidence confirms this mismatch. As shown in \textbf{Tab.}~\ref{tab:CV8DResults}, CV-8D exhibits many misses (avg 645) and a high counting error (MAPE 14.59\%). The root cause is that during occlusions, the constant-velocity Kalman filter cannot predict the object’s reappearance position and size accurately. Large deviations between prediction and observation (i.e., large Mahalanobis distance) lead to association failures: the tracker either terminates the original trajectory (misses) or allocates a new identity upon reappearance (ID switches). The performance bottleneck is thus that constant velocity assumptions cannot effectively predict accurate positions in the next frame.

\subsubsection{Gains and Losses of the Higher-Order Kinematic Model (CA-12D)}

A straightforward improvement to CV-8D is to raise the kinematic order by adding accelerations. The CA-12D model embodies this idea and can, in principle, fit nonlinear motion trajectories.
The data partially validate this: relative to CV-8D, CA-12D improves MOTA from 64.11\% to 67.54\% and reduces counting MAPE from 14.59\% to 6.99\%. This stems mainly from better prediction after short occlusions, which increases successful associations and substantially reduces false positives (avg from 480 to 278).
However, CA-12D introduces an unintended drawback: identity switches rise from 24 to 39.15, even exceeding those of the simpler CV-8D. This reveals a typical challenge in high-dimensional state estimation: while higher dimensionality increases expressiveness, it also amplifies sensitivity to measurement noise. YOLOv11 detections inevitably contain small frame-to-frame jitters. CA-12D may interpret such tiny jitters as large instantaneous accelerations, destabilizing state estimates. Unstable predictions disrupt association, particularly in crowded, ambiguous settings, and increase identity mis-assignments. The lesson is clear: without strong constraints, simply adding degrees of freedom may improve trajectory fitting but harm robustness to noise and, consequently, identity consistency.

\subsubsection{Why the Physics-Constrained Model (Phys-5D) Prevails}

Motivated by CA-12D’s trade-offs, we propose Phys-5D, which lifts estimation from the 2D image plane to the 3D physical domain and leverages inherent geometric constraints.
The key is to use the pinhole camera model to establish a deterministic, nonlinear relation between the 2D image height $h$ (px) and the 3D depth $Z$. This acts as a strong regularizer: height jitters are no longer arbitrary size changes but physically interpreted as small movements along depth. Advantages include: (1) physical realism by enforcing perspective consistency and suppressing artifacts; (2) incorporation of priors via an adaptive deceleration constraint that reflects real behavior as pedestrians approach the camera/train; and (3) reduced state dimensionality and improved stability—a compact 5D state with a highly constrained solution space enables robust inference from noisy observations, directly addressing CA-12D’s instability.
Empirically, Phys-5D achieves the best balance (\ref{tab:Phys5DResults}). Its MOTA (67.21\%) is comparable to CA-12D, but ID switches regress to baseline levels (24.80, on par with CV-8D). This stability translates to the application level: counting MAPE drops to 2.24\%.
To visualize the trade-offs and evolution across models, we summarize key metrics below.
As in Table \ref{tab:3modelsCountingComparisonResults}, a comparison of the three motion models:

\begin{table}[H]
    \centering
    \resizebox{0.6\textwidth}{!}{%
    \small
    \begin{tabular}{ccccc}
    \hline
        \textbf{Model} & \textbf{MAE } & \textbf{RMSE} & \textbf{MAPE(\%)} & \textbf{ME } \\ \hline
        \textbf{CV-8D} & 3.4 & 6.1514 & 14.592 & 1.4  \\ 
        \textbf{CA-12D} & 2.4 & 4.5837 & 6.99 & 0.7  \\ 
        \textbf{Phys-5D} & 0.75 & 1.3229 & 2.236 & -0.25 \\ \hline
    \end{tabular}%
    }
    \caption{Counting Performance Comparison of the Three Models}
    \label{tab:3modelsCountingComparisonResults}
\end{table}

The table makes the evolution clear: from CV-8D to CA-12D, we improved motion fitting at the expense of identity stability; from CA-12D to Phys-5D, imposing physics constraints not only restored and surpassed identity stability but also pushed application accuracy to a new level. Phys-5D is thus the optimal solution for our task.

\subsection{Synergy and Contributions of System Components}

Although the motion model is central, the system’s success results from carefully designed components working together. The detect–track–count pipeline is a chain whose overall strength depends on every link.
Detector as the foundation: the upper bound of tracking is largely determined by detector quality. Our two-stage strategy (generic pre-training + domain fine-tuning) raises mAP50–95 from 54.7\% to 81.6\%. High-precision, well-localized detections provide low-noise measurements to the Kalman filter—crucial for all motion models, especially noise-sensitive CA-12D and geometry-reliant Phys-5D. Without a strong, scene-optimized detector, sophisticated motion modeling would be built on sand.
Re-ID as the safety net: DeepSORT relies on both motion and appearance cues. We train a lightweight, high-performance EfficientNet-B0 Re-ID specifically for head images in platform scenes, achieving 100\% Rank-1 on our test set. Robust appearance descriptors serve as a safety net when motion prediction falters under long occlusions or abrupt movements, enabling correct re-association and preventing ID switches. This complementarity underpins the tracker’s robustness.
Counting strategy as the final mile: even a perfect tracker cannot yield accurate counts with a brittle counting logic. The ablations in Chapter 5 show that a single counting line fails catastrophically (MAPE 93.43\%) even with the best tracking input, due to extreme sensitivity to jitter and brief detection gaps. Our virtual counting band, requiring persistence across frames within a nonzero-width zone and applying end-of-video compensation, dramatically improves robustness and converts high-quality tracks into accurate counts.
In short, our success reflects both divide-and-conquer and co-optimization: by optimizing detection, re-identification, motion prediction, and counting logic and then integrating them cohesively, we achieve system-level excellence.

\subsection{Theoretical Implications and Practical Value}

Theoretical implications: Our core contribution is a compelling case for Physics-Informed Machine Learning (PIML) in computer vision. While recent trends often favor bigger, deeper, purely data-driven models, our results show that in scenes governed by clear physical laws, injecting first principles (e.g., perspective geometry) can yield larger gains and greater robustness than merely adding parameters (e.g., 8D→12D). This challenges brute-force data-driven thinking and advocates physics-aware modeling.
Practical value: Our results directly address the real-world need raised in Chapter 1. A 2.24\% counting error is deployment-ready. For railway operators, accurate counts enable: \\(1) Safety management: real-time density monitoring with automatic alerts near thresholds to prevent overcrowding and accidents; \\(2) Operations optimization: precise knowledge of temporal and per-train passenger patterns for schedule planning and dynamic frequency adjustment, improving efficiency and reducing waiting time; \\(3) Emergency response: reliable counts for evacuation planning during disruptions (e.g., delays). Thus, beyond academic value, our approach is a practical solution that can enhance safety, efficiency, and passenger experience in railway systems.

\subsection{Limitations}

Despite promising results, we critically examine limitations to guide future work. Camera-parameter dependence: Phys-5D’s strength—accurate camera geometry—also limits plug-and-play deployment, as intrinsics (focal length, principal point) are required. While lab calibration is straightforward, large-scale manual calibration is impractical.
Scene/model specificity: Phys-5D is tailored to a camera mounted on a moving train observing a static platform, with priors (e.g., adaptive deceleration) tied to this setup. Direct deployment to other settings (e.g., fixed overhead cameras, strong camera shake) may degrade performance, reflecting limited out-of-domain generalization.
Dataset coverage: Our MOT-RPCH is constructed primarily from public videos and may lack extreme conditions (night, rain/snow, harsh lighting), leaving generalization to out-of-distribution scenarios insufficiently validated.
Simplified physics: For efficiency, we fix head physical height H at 0.3 m and consider only depth-direction motion, omitting full 6-DoF head pose. While effective here, such simplifications may bottleneck tasks requiring finer pose estimation.

\iffalse

本章对第5章呈现的实验结果提供全面的分析和解释。我们调查三种卡尔曼状态空间配置之间性能差异的根本原因，阐明提出的物理约束模型（Phys-5D）如何实现卓越性能，并分析系统组件之间的协同相互作用。我们进一步讨论理论意义、实践价值和局限性，为未来研究方向提供见解。

6.1 结果的综合解读与核心发现

第5章的实验结果证明了我们检测-跟踪-计数框架的有效性。我们的最终系统实现了最先进的计数精度，MAPE=2.24%，相比基线方法有显著改进。关键发现来自我们全面评估：

\textbf{性能总结：}三种卡尔曼状态空间配置显示出不同的性能特征：
\begin{itemize}
    \item CV-8D（基线）：MAPE=14.59\%，MAE=3.4，ID切换=24
    \item CA-12D（加速度）：MAPE=6.99\%，MAE=2.4，ID切换=39.15
    \item Phys-5D（物理约束）：MAPE=2.24\%，MAE=0.75，ID切换=24.8
\end{itemize}

\textbf{核心发现：}结果表明，将领域知识（相机几何和物理约束）融入运动模型比仅仅增加运动学模型复杂度更有效。Phys-5D在跟踪精度和身份一致性之间实现了最佳平衡，转化为卓越的计数性能。

本次研究的最终成果：一个集成了经过领域自适应微调的YOLOv11m检测器、定制化训练的EfficientNet-B0重识别（Re-ID）模型以及创新的5维物理约束卡尔曼滤波器（Phys-5D）的跟踪计数系统在关键的应用层指标上取得了卓越的性能。在自建的MOT-RPCH测试集上，该系统的平均绝对百分比误差（MAPE）仅为2.24\%。这一高精度的计数结果，与作为基准的标准8维匀速模型（CV-8D）高达14.59\%的MAPE，以及引入加速度项的12维模型（CA-12D）的6.99\%的MAPE形成了鲜明对比。这一性能层级上的显著差异，构成了本章后续深入讨论的基石。

一个尤为深刻的发现是，最终的应用性能（计数精度）与传统的多目标跟踪（MOT）指标（如MOTA）之间并非简单的线性正相关关系。实验数据显示，CA-12D模型与Phys-5D模型取得了几乎相当的MOTA分数（分别为67.54\%和67.21\%），这表明两者在逐帧检测与关联的综合准确率上表现相近。然而，若仅凭此指标，则会忽略一个关键差异：CA-12D模型的身份切换（ID Switches）次数（平均39.15次）远高于Phys-5D模型（平均24.80次）。ID切换次数是衡量跟踪器维持身份一致性能力的核心指标。对于计数这类需要对个体进行长期、稳定、唯一身份识别的累积性任务而言，轨迹的稳定性和身份的持久性远比单帧的定位精度更为重要。每一次身份切换都可能导致系统将同一个体误判为两个不同的个体，从而引发重复计数或漏计。因此，尽管CA-12D在拟合瞬时运动上表现尚可，但其身份保持能力的欠缺直接导致了其计数误差远高于Phys-5D。这一现象揭示了，在面向应用的跟踪系统评估中，必须更加重视如IDF1和ID切换次数这类衡量身份一致性的指标。Phys-5D模型的真正优势不仅在于准确地跟踪目标，更在于其能够持续、稳定地识别目标。

6.2 运动模型：从纯运动学到物理约束的演进与权衡

为了系统性地揭示不同运动模型性能差异的根源，本节将对实验中采用的三种模型——CV-8D、CA-12D和Phys-5D——进行细致的比较分析，论证为何引入物理约束是实现本次任务成功的关键。

6.2.1 基准模型 (CV-8D) 的局限性分析

作为多目标跟踪领域的标准基准，8维匀速模型（CV-8D）构成了我们评估的起点。该模型的基本假设是目标在图像平面内以恒定速度进行线性运动。然而，在我们的特定场景——列车视角下的站台人群——这一假设被从根本上打破。随着列车进站减速，由于透视投影效应，站台上的行人在图像平面内的视在运动呈现出显著的非线性加速特征，其边界框尺寸也随之非线性地放大。
实验结果有力地印证了这一理论上的不匹配。如\ref{tab:CV8DResults}所示，CV-8D模型表现出大量的漏检（Misses，平均645个），最终导致了高达14.59\%的计数误差。其失败的根源在于，当目标因人群遮挡而短暂消失时，基于匀速假设的卡尔曼滤波器无法准确预测其重现时的位置和尺寸。预测与实际观测之间的巨大偏差（即马氏距离过大）导致数据关联失败，跟踪器要么错误地终止原有轨迹（造成漏检），要么在目标重现时为其分配一个新的身份（造成ID切换）。因此，CV-8D模型的性能瓶颈，是其运动学假设与场景物理现实之间的根本性矛盾。

6.2.2 高阶运动学模型 (CA-12D) 的得与失

为了克服匀速模型的局限性，一个直观的改进思路是提升模型的运动学阶数，即引入加速度项。12维恒定加速度模型（CA-12D）正是基于这一理念设计的。通过在其状态向量中增加加速度分量，该模型在理论上具备了拟合非线性运动轨迹的能力。
实验数据在一定程度上验证了这一改进的有效性。相较于CV-8D，CA-12D模型的MOTA从64.11\%提升至67.54\%，计数MAPE也从14.59\%显著降低至6.99\%。这主要归功于其更强的轨迹预测能力，尤其是在短时遮挡后，更准确的预测位置提高了关联成功的概率，从而大幅减少了误报（False Positives）数量（从平均480个降至278个）。
然而，CA-12D模型也带来了一个意想不到的负面效应：其身份切换次数不降反升，从24次急剧增加到39.15次，甚至超过了更简单的CV-8D模型。这一现象揭示了高维状态空间估计中的一个典型困境,即更高维度的模型虽然表达能力更强，但同时也对输入的测量噪声更为敏感。前端的YOLOv11检测器所输出的边界框位置不可避免地存在微小的、逐帧的抖动（即测量噪声）。对于CA-12D模型而言，卡尔曼滤波器会试图从这些带有噪声的位置信息中估计出加速度。一个微小的位置抖动，就可能被错误地解释为一个剧烈的、瞬时的加速度变化，从而导致状态估计变得极不稳定。一个不稳定的状态预测会严重干扰数据关联过程，尤其是在人群密集、多个目标相互靠近的模糊情境下，更容易导致身份的错误指派。因此，CA-12D模型的“得与失”为我们提供了一个深刻的教训：在没有强力约束的情况下，单纯增加模型的自由度（维度），虽然可能提升其对理想轨迹的拟合能力，但代价是牺牲了对现实世界中噪声的鲁棒性，反而可能损害对身份一致性至关重要的跟踪稳定性。

6.2.3 物理约束模型 (Phys-5D) 的优越性论证

正是在对CA-12D模型权衡的反思之上，本研究提出了创新的5维物理约束模型（Phys-5D）。该模型的核心思想是，与其在2D图像平面上构建一个日益复杂的纯运动学模型，不如将问题提升至3D物理空间，并利用场景固有的几何约束来规范化状态估计过程。
Phys-5D模型的关键在于利用针孔相机模型，在目标的2D图像高度 h（px）​ 与其在3D空间中的物理深度Z之间建立了一个确定性的、非线性的强约束关系。这一约束关系扮演了一个强大的正则化器（regularizer）角色。它从根本上改变了状态估计的方式：检测框高度的抖动不再被解释为物体尺寸的任意变化，而是被严格地、物理地解释为物体在深度方向上的微小位移。这种做法带来了多重优势：
（1）物理真实性：它确保了目标的尺寸变化严格遵循透视几何规律，有效抑制了伪影。
（2）先验知识融入：模型中引入的自适应减速度约束，模拟了行人在列车靠近时通常会减速的真实行为模式，进一步增强了预测的合理性。
（3）状态空间降维与稳定：尽管模型背后的物理原理更复杂，但其状态向量仅有5维，比CA-12D更为紧凑。更重要的是，物理约束极大地压缩了状态估计的解空间，使得滤波器能够从带噪的观测中稳健地推断出物理上一致的状态，从而直接解决了CA-12D模型因过拟合噪声而导致的状态不稳定问题。

实验结果全面证实了Phys-5D模型的优越性。如\ref{tab:Phys5DResults}所示，该模型在各项指标上取得了最佳的平衡。其MOTA（67.21\%）与CA-12D相当，但ID切换次数（24.80次）成功地降回了与最简单的CV-8D模型持平的水平。这种在保持高精度的同时兼顾身份稳定性的能力，最终转化为应用层面的巨大成功：计数MAPE降至仅2.24\%。
为了更直观地展示三者之间的权衡与演进，下表总结了关键的性能指标。
如表 \ref{tab:3modelsCountingComparisonResultsCN}所示，三种运动模型关键性能指标对比

\begin{table}[!ht]
    \centering
    \resizebox{0.4\textwidth}{!}{%
    \small
    \begin{tabular}{ccccc}
    \hline
        \textbf{Model} & \textbf{MAE } & \textbf{RMSE} & \textbf{MAPE(\%)} & \textbf{ME } \\ \hline
        \textbf{CV-8D} & 3.4 & 6.1514 & 14.592 & 1.4  \\ 
        \textbf{CA-12D} & 2.4 & 4.5837 & 6.99 & 0.7  \\ 
        \textbf{Phys-5D} & 0.75 & 1.3229 & 2.236 & -0.25 \\ \hline
    \end{tabular}%
    }
    \caption{三种运动模型计数性能对比}
    \label{tab:3modelsCountingComparisonResultsCN}
\end{table}

该表格清晰地揭示了研究的演进逻辑：从CV-8D到CA-12D，我们以牺牲身份稳定性为代价换取了运动拟合能力的提升；而从CA-12D到Phys-5D，我们通过引入物理约束，不仅恢复并超越了原有的身份稳定性，还进一步将应用精度推向了新的高度。这强有力地证明，Phys-5D是本次研究任务的最优解决方案。

6.3 系统各组件的协同效应与贡献

尽管创新的运动模型是本研究的核心，但最终系统取得的成功，是其各个精心设计的组件协同作用的结果。整个“检测-跟踪-计数”流程如同一条环环相扣的链条，任何一个环节的薄弱都会导致最终性能的下降。

检测器是基石：跟踪系统的性能上限，很大程度上由前端检测器的质量决定。本研究采用的“通用预训练 + 领域微调”两阶段策略，成功地将YOLOv11m检测器在关键的定位精度指标mAP50-95上从54.7\%提升至81.6\% 。一个高精度、高定位准确率的检测器，能够为后续的卡尔曼滤波器提供低噪声的测量输入。这对于所有运动模型，尤其是对噪声敏感的CA-12D和依赖精确几何关系的Phys-5D而言，是其能够稳定运行的根本前提。可以说，没有一个强大的、为特定场景优化的检测器，后续关于运动模型的精妙设计都将是空中楼阁。

Re-ID模型是保障：DeepSORT框架的精髓在于运动信息与外观信息的结合。在本研究中，我们并未使用通用的Re-ID模型，而是针对站台场景中的头部图像，专门训练了一个基于EfficientNet-B0的轻量级、高性能模型，其在我们的测试集上达到了100\%的Rank-1准确率。这个强大的外观特征提取器，在跟踪系统中扮演着“安全网”的角色。当运动模型因长时间遮挡或剧烈运动而预测失效时，鲁棒的外观特征能够成为数据关联的最后一道防线，准确地将重现的目标与其原始身份重新关联，从而有效防止ID切换。运动模型与Re-ID模型的这种互补协同，是整个跟踪器鲁棒性的重要保障。

计数策略是终点：一个完美的跟踪器如果匹配了拙劣的计数逻辑，同样无法产出准确的结果。第五章中关于计数方法的消融实验戏剧性地证明了这一点：传统的单计数线方法，即便输入的是我们最优的跟踪结果，其产生的MAPE也高达93.43\%，几乎完全失效。这揭示了计数线方法对轨迹的微小抖动和检测的瞬时中断极其敏感，存在严重的系统性漏计偏差。我们为此设计的虚拟计数带策略，通过要求目标在一定宽度的区域内持续出现数帧才进行计数，并辅以最后帧补偿机制，极大地增强了计数过程对噪声的鲁棒性。这一策略的有效性，是我们将高质量的跟踪轨迹转化为高精度计数结果的“最后一公里”。

综上所述，本研究的成功并非源于单一的技术突破，而是体现了系统工程中“分而治之”与“协同优化”的思想。通过对检测、重识别、运动预测和计数逻辑这四个核心模块的逐一优化与精心集成，我们构建了一个各部分性能相互增强、整体表现卓越的系统。

6.4 研究的理论意义与实践价值

理论意义：本研究最核心的理论贡献，是为计算机视觉领域内日益受到关注的“物理信息机器学习”（Physics-Informed Machine Learning, PIML）范式提供了一个强有力的成功案例。在深度学习浪潮下，许多研究倾向于使用更大、更深、更复杂的纯数据驱动模型来解决问题。而我们的实验清晰地表明，在一个具有明确物理规律的场景中，将第一性原理（如透视几何）作为强先验融入模型，能够比单纯增加模型参数（如从8维增至12维）带来更显著的性能提升和更强的鲁棒性。这项工作挑战了“暴力美学”的数据驱动思想，倡导了一种数据与物理知识相结合的建模新思路，对于其他具有类似物理约束的跟踪任务（如基于阿克曼转向几何的车辆跟踪、基于流体动力学的粒子跟踪等）具有重要的启发意义。

实践价值：本研究的成果具有直接且显著的实践价值，直接回应了第一章中提出的现实需求。2.24\%的计数误差率，意味着系统已经达到了可以投入实际应用的精度水平。对于铁路运营商而言，如此精准的客流数据是实现智能化运营的关键：
（1）安全管理：实时监测站台的人群密度，当人数接近或超过安全阈值时自动预警，有效预防拥挤、踩踏等安全事故。
（2）运营优化：通过分析历史客流数据，运营商可以精确掌握不同时段、不同车次的客流模式，从而优化列车排班、动态调整发车频率，提升运输效率，减少乘客等待时间。
（3）应急响应：在突发事件（如列车延误）导致乘客滞留时，精确的计数系统能够为应急疏散方案的制定提供关键的数据支持。
因此，本研究不仅在学术上探索了新的模型范式，更提供了一套具备高应用价值的技术解决方案，有望为提升铁路运输系统的安全性、效率和乘客体验做出实质性贡献。

6.5 研究的局限性

尽管本研究取得了令人鼓舞的成果，但我们必须以批判性的眼光审视其存在的局限性，这既是严谨学术态度的体现，也为未来的研究指明了方向。
相机参数的依赖性：Phys-5D模型的最大优势来源于其对相机几何的精确建模，但这同时也构成了其最大的局限。该模型需要预先知道相机的内参（焦距、主点）。在实验室环境下，这可以通过标准的棋盘格标定法获取。但在大规模实际部署中，为成百上千个摄像头逐一进行手动标定是不现实的，这极大地限制了该模型的“即插即用”能力。
场景与运动假设的特化性：本研究中的物理模型是为“固定于移动列车上的相机观测静态站台”这一高度特定的场景量身定制的。模型中的自适应减速度等先验假设，都与此场景紧密相关。如果将该模型直接应用于其他场景，如安装在站台天花板上的固定摄像头，或者存在显著相机抖动的情况，其性能可能会大幅下降。模型的特化性在带来高精度的同时，也牺牲了一定的泛化能力。
数据集的局限性：我们自建的MOT-RPCH数据集是本次研究能够成功进行的关键 。然而，该数据集主要来源于公开视频，尽管我们尽力使其多样化，但其在覆盖所有真实世界变化方面仍有不足，例如，缺乏夜间、雨雪等恶劣天气和极端光照条件下的数据。因此，模型在这些未曾见过的、分布外（out-of-distribution）场景下的泛化性能尚未得到充分验证。
简化的物理模型：为了模型的简洁与高效，我们在物理建模中做出了一些简化假设。例如，我们假设所有行人的头部物理高度 H 均为一个固定值（0.3米），这与现实中身高差异不符。此外，模型只考虑了深度方向的运动，并未对头部的完整六自由度（6-DoF）姿态进行建模。虽然这些简化在当前任务中被证明是有效的，但在需要更精细姿态估计的任务中可能会成为瓶颈。
\fi